\documentclass{beamer}
% Load Packages
\usepackage[utf8]{inputenc}
\usepackage{xcolor}
\usepackage{tikz}
\usetikzlibrary{positioning,calc}
\usepackage{graphicx}
\usepackage{hyperref}
\usepackage{amsmath}
\usepackage{listings}
\usepackage{fontawesome}
\usepackage{hyperref}

% Define Commands
\newcommand*{\ClipSep}{0.06cm} %To adjust footer logo
% \newcommand{\E}{\mathrm{e}\,} %\def\I{e} % used to defined e for exp(x), see later what it should be
\newcommand{\ud}{\mathrm{d}}
\newcommand{\heart}{\ensuremath\heartsuit}
\DeclareMathOperator{\E}{\mathbb{E}}
\lstset{numbers=left, numberstyle=\tiny, stepnumber=1,firstnumber=1,breaklines=true,
    numbersep=5pt,language=Python,
    stringstyle=\ttfamily,
    basicstyle=\footnotesize, 
    showstringspaces=false
}

\setlength{\parskip}{\baselineskip}

\usetheme{oxonian}

%--------------------------- Title Slide

\title{Accelerating Entropy-Based Transformer Calibration}
\titlegraphic{\includegraphics[width=3cm]{Theme/Logos/pton-shield.png}}
\author{Michael Hu}
\institute{Advisor: Karthik Narasimhan}
\date{} %\today

\begin{document}

{\setbeamertemplate{footline}{} 
\frame{\titlepage}}

%---------------------------- Outline

% \section*{Outline}
% \begin{frame}{Outline}
% \tableofcontents
% \end{frame}

%---------------------------- Motivation

\section{Motivation}

\begin{frame}{Motivation: Preliminaries}

\begin{enumerate}
    \item Language model: $\widehat{Pr}(w_t | w_{<t})$
    \pause
    \item Entropy: a measure of the amount of information being produced by a source.
    \pause
    \item Entropy rate: the average amount of information per word. For a sentence, sum the entropies and divide by number of words.
    \pause
        \begin{itemize}
            \item "A loud and boisterous carnival is outside my window."  
            $\Rightarrow$ High entropy rate.
            \pause
            \item "Nothing is outside my window." $\Rightarrow$ Low entropy rate.
            \pause
        \end{itemize}
    \item Entropy blowup: when the entropy rate of text produced by a source increases over time.
\end{enumerate}
\end{frame}

\begin{frame}{Motivation}
Entropy blowup does not occur in human-generated text.
\begin{itemize}
    \item Intuition: people don't say increasingly weird things as you talk to them.
\end{itemize}
\pause
We observe entropy blowup on SOTA language models such as GPT-2. 
\end{frame}

\section{Goal}
\begin{frame}{Goal}
Develop a computationally efficient technique for minimizing entropy blowup in language models.
\end{frame}

\section{Related Work}
\begin{frame}{Related Work}
    \begin{itemize}
        \item Analysis of language model sampling methods by Holtzman et al. \cite{holtzman2019curious}
        \item Calibration of neural networks via temperature scaling by Guo et al. \cite{guo2017calibration}
    \end{itemize}
\end{frame}

\section{Approach}
\begin{frame}{Approach: Preliminaries}
Conditional entropy: 
\begin{align*}
    H(w_t|w_{<t}) := \E_{w_t \sim Pr} \left[ \log \frac{1}{Pr(w_t|w_{<t})} \right]
\end{align*}
\pause
Conditional one-step lookahead entropy:
\begin{align*}
    H(\widehat{W}_{t+1} | w_{\leq t}) := \E_{w_{t+1} \sim \widehat{Pr}} \left[\log \frac{1}{\widehat{Pr}(w_{t+1} | w_{\leq t} )} \right]
\end{align*}
\end{frame}

\begin{frame}{Approach}
    Stepwise calibration:
    \begin{enumerate}
        \item Obtain the learned language model $\widehat{Pr}$.
        \pause
        \item Scale $\widehat{Pr}$ by $\exp\{-\alpha \cdot H(\widehat{W}_{t+1} | w_{\leq t})\}$. Renormalize to obtain a new probability distribution $Pr_\alpha$.
        \pause
        \item Choose $\alpha$ such that it minimizes the cross-entropy loss between the context and $Pr_\alpha$.
        \pause
    \end{enumerate}
    Use $Pr_\alpha$ as your new language model.
\end{frame}

\begin{frame}{Approach}
    \begin{center}
        \includegraphics[width=8cm]{figures/calibration_effect.png}
    \end{center}
\end{frame}

\begin{frame}{Approach: How to improve?}
Computational complexity:
\begin{itemize}
    \item V = size of vocabulary = $\sim$ 50,000 in practice
    \item C = number of contexts = for a sentence, the number of words $\Rightarrow$ $\sim$ 10. 
    \item T = number of words we wish to generate
\end{itemize}
\pause
The cross-entropy loss takes $O(V^2C)$ time to compute. 

Computing $Pr_\alpha$ takes $O(V^2)$ time $\Rightarrow$ sampling $T$ words takes $O(V^2T)$ time.
\end{frame}

\section{Implementation}
\begin{frame}{Implementation}
    The major contribution of this IW: we can approximate the true one step lookahead entropy by calculating lookahead entropies for the top 100 or so words in the vocabulary.
    \pause

    Impact:
    \begin{itemize}
        \item $V$= 50,000 $\Rightarrow$ 100. If the model used to take 8 hours to calibrate or generate $T$ words, it now takes 1 minute.
    \end{itemize}
\end{frame}

\section{Results}
\begin{frame}{Result 1}
    On average, top 128 words capture $>90\%$ of the probability mass.
    \begin{center}
        \includegraphics[width=6cm]{figures/top128_pm.png}
    \end{center}
    Here: 92.8\%
\end{frame}

\begin{frame}{Result 2}
    Generating using approximations of $Pr_\alpha$ minimizes entropy blowup.
    \begin{center}
        \includegraphics[width=8cm]{figures/different_calibrations.png}
    \end{center}
\end{frame}

\begin{frame}{Result 3}
    We observe no qualitative drawbacks when using approximations of $Pr_\alpha$.
    % TODO: add generations, show that they are good.
\end{frame}

\begin{frame}{Result 4}
    A plug-and-play API for calibration, available at \url{https://github.com/mikkyhu/transformers/blob/master/examples/calibrate.py}
\end{frame}

\section{Conclusion}
\begin{frame}{Conclusion}
    Approximating $Pr_\alpha$ is a viable way to speed up calibration and generating from $Pr_\alpha$.
\end{frame}

\section{Future Work}
\begin{frame}{Future work}
    \begin{center}
        \includegraphics[width=8cm]{figures/different_calibrations_top40.png}
    \end{center}
\end{frame}

\section{Acknowledgements}
\begin{frame}{Acknowledgements}
    \begin{itemize}
        \item Karthik Narasimhan
        \item Cyril Zhang
    \end{itemize}
    \bibliographystyle{unsrt}
    \bibliography{References}
\end{frame}

\end{document}

